\documentclass{article}
\usepackage{amsmath}
\begin{document}
\section{Offline Reference Support at Joins}
The existing computed-torque tracker linearly interpolates position and
interpolates velocity and acceleration tables formed by applying
\texttt{numpy.gradient} twice. For an interior uniform grid away from
segment boundaries, the stencils are
\[
v_i=\frac{q_{i+1}-q_{i-1}}{2h},\qquad
a_i=\frac{q_{i+2}-2q_i+q_{i-2}}{4h^2}.
\]
At endpoints the library uses first-order one-sided differences; nonuniform
grids use different weights. An off-grid interpolation may use two adjacent
rate rows, extending the set of source indices. The diagnostic returns a
conservative structural union, including terms that may cancel, in the input
clock supplied by the caller. Invalid or numerically overflowing clocks fail.

GCV-20 differentiates and interpolates rate tables separately on each side
of impact. The existing position sample is unsplit, so a query
before the split can use a position row after it. This describes existing
offline behavior and does not establish an implementation defect. The opt-in
past-only guard rejects any structurally possible row later than the query.
Neither the guard nor a passing result audits upstream filtering, target
generation, reference hooks or closed-loop dynamics; no physical or causal
inference follows. Test by perturbation and basis vectors with
\texttt{python -m pytest -q
tests/unit/motion\_matching/test\_reference\_support\_15493.py}.
\end{document}
